\begin{textsample}{NEG Dim 1 – 2017\_c5 – Score 14.00 – 2017\_c5/enade\_2017\_letras\_por\_281c81746d.txt}  \label{ex:f1_neg_390}
CAPÍTULO 5 DISTRIBUIÇÃO DOS CONCEITOS 5.1 PANORAMA NACIONAL DA DISTRIBUIÇÃO DOS CONCEITOS O Conceito Enade é um indicador de qualidade que avalia os cursos por intermédio dos desempenhos dos estudantes a partir dos resultados do Enade.
Os valores possíveis vão de 1 ( pior situação ) a 5 ( melhor situação ).
Os cursos com apenas um ou sem nenhum concluinte participante não obtêm o Conceito Enade, ficando Sem Conceito ( SC ).
Na Tabela 5.1a, são apresentadas a quantidade e a distribuição de cursos de Licenciatura de Letras-Português participantes do Enade/2017, por faixa de conceito e Grande Região.
Enfatiza -se, mais uma vez, que a diferença dentre os cursos tabulados neste capítulo e no capítulo 2 corresponde aos cursos Sem Conceito ( SC ), ou seja, àqueles, em princípio, sem alunos concluintes que houvessem participado da prova ou àqueles que tivessem tão somente um aluno concluinte18.
Observando -se os dados da Tabela 5.1a, nota -se que, dos 318 cursos participantes, 108 ( 34,0 \% ) classificaram -se com conceito 3, o valor modal.
Esse foi também o conceito modal nas regiões Nordeste ( 41,6 \% ) e Centro-Oeste ( 41,9 \% ).
Nas regiões Sudeste ( 33,3 \% ) e Sul ( 43,5 \% ), o conceito modal foi 4.
Já na região Norte, a moda foi o conceito 2, atribuído a 34,2 \% dos cursos da região.
O conceito 4 foi o segundo mais frequente em nível nacional ( 24,5 \%, correspondendo a 78 cursos ) e o conceito 2, o terceiro ( 23,0 \%, correspondendo a 73 cursos ).
Houve, ainda, 20 cursos ( 6,3 \% ) que receberam conceito 5, e 24 cursos ( 7,5 \% ) que receberam conceito 1.
Dos 318 cursos de Letras-Português, 15 ( 4,7 \% ) ficaram Sem Conceito ( SC ). 18 Esses cursos com somente um concluinte não foram considerados no capítulo 2.
A região Norte participou com 38 cursos ou 11,9 \% do total nacional.
Desses, 13 cursos receberam o conceito 2, o que equivale a 34,2 \% do total regional, além de ser o conceito modal na região, como já comentado.
A seis cursos ( 15,8 \% ) atribuiu -se conceito 1, a 12 cursos ( 31,6 \% ), conceito 3, e a cinco cursos ( 13,2 \% ), conceito 4.
Dois cursos ( 5,3 \% ) ficaram Sem Conceito ( SC ), e nenhum curso recebeu conceito 5 nessa região.
A região Nordeste participou com 113 cursos ou 35,5 \% do total nacional.
Nessa região, 47 cursos ( 41,6 \% em termos regionais ) obtiveram conceito 3, o conceito modal para a região, como já comentado.
Os conceitos 1 e 2 foram atribuídos, respectivamente, a 11 e a 32 cursos ( 9,7 \% e 28,3 \%, respectivamente ).
Ainda nessa região, 19 cursos ( 16,8 \% ) foram avaliados com o conceito 4, e dois cursos ( 1,8 \% ), com conceito 5.
Dois cursos ( 1,8 \% ) ficaram Sem Conceito ( SC ).
Dos 90 cursos participantes da região Sudeste ( 28,3 \% do total nacional ), o conceito 4 foi o conceito modal, como já anteriormente assinalado, tendo sido obtido por 30 cursos ( 33,3 \% ).
O conceito 1 foi atribuído a dois cursos ( 2,2 \% ) e o conceito 2, a 17 cursos ( 18,9 \% ).
Receberam os conceitos 3 e 5, respectivamente, 21 cursos ( 23,3 \% ) e 12 cursos ( 13,3 \% ).
Nessa região, oito cursos ( 8,9 \% ) ficaram Sem Conceito ( SC ).
Os 46 cursos da região Sul corresponderam a 14,5 \% do total nacional.
Como já indicado, houve predominância do conceito 4, atribuído a 43,5 \% dos cursos da região ( conceito modal ), ou seja, atribuído a 20 dos 46 cursos participantes na região Sul.
O conceito 1 não foi atribuído a nenhum curso, o conceito 2 foi atribuído a cinco cursos ( 10,9 \% ), o conceito 3, a 15 cursos ( 32,6 \% ), e o conceito 5, a cinco cursos ( 10,9 \% ).
Um curso ( 2,2 \% ) ficou Sem Conceito ( SC ) na região Sul.
Já dos 31 cursos participantes na região Centro-Oeste ( 9,7 \% do total nacional ), 13 receberam conceito 3, o conceito modal ( 41,9 \% ), como já fora destacado.
A cinco cursos ( 16,1 \% ) foi atribuído conceito 1, e a seis cursos ( 19,4 \% ), conceito 2.
Os conceitos restantes, 4 e 5, foram atribuídos, respectivamente, a quatro cursos ( 12,9 \% ) e a um curso ( 3,2 \% ).
Nessa região, dois cursos ( 6,5 \% ) ficaram Sem Conceito ( SC ).
No Gráfico 5.1a, apresenta -se a distribuição cumulativa do Conceito Enade dos cursos de Licenciatura, segundo a Grande Região.
Em um gráfico de distribuição cumulativa, em \textbf{cada} valor de abscissa ( eixo x ), o valor apresentado no eixo das ordenadas ( eixo y ) é o valor acumulado da variável.
Por exemplo, a região Norte ( linha azul ) apresenta um pouco mais de 86 \% dos cursos com conceito 3 ou abaixo deste valor.
A região Sul ( linha roxa ), por outro lado, apresenta um pouco mais de 44 \% dos cursos na mesma situação : com conceito 3 ou abaixo.
Na região Norte ( linha azul ), todos os cursos têm conceitos 4 ou abaixo disso, já que a poligonal cumulativa alcança 100 \% para o conceito 4.
Em linhas gerais, em um gráfico de distribuição cumulativa, poligonais mais à esquerda representam Grandes Regiões com pior distribuição, e poligonais mais à direita, Grandes Regiões com melhores conceitos.
A região Norte apresenta os cursos com conceitos mais baixos, ao passo que as regiões Sul ( linha roxa ) e Sudeste ( linha vermelha ) apresentam as distribuições com valores maiores.
Cursos em IES situadas nas regiões Nordeste ( linha verde ) e Centro-Oeste ( linha preta ) \textbf{estão} em uma situação intermediária, porém mais próxima da região Norte ( linha azul ) do que as ( linha roxa ) e Sudeste ( linha vermelha ).
Na Tabela 5.1b, são apresentadas a quantidade e a distribuição de cursos de Bacharelado de Letras-Português participantes do Enade/2017, por faixa de conceito e Grande Região.
Enfatiza -se, mais uma vez, que a diferença dentre os cursos tabulados neste capítulo e no capítulo 2 corresponde aos cursos Sem Conceito, ou seja, àqueles, em princípio, sem alunos concluintes que houvessem participado da prova ou àqueles que tivessem tão somente um aluno concluinte¹⁹. ¹⁹ Estes cursos com somente um concluinte não foram considerados no capítulo 2.
Observando -se os dados da Tabela 5.1b, nota -se que, dos 34 cursos participantes, 16 ( 47,1 \% ) classificaram -se com conceito 3, o valor modal.
Esse foi também o conceito modal nas regiões Nordeste ( 100,0 \% ) e Sudeste ( 50,0 \% ).
Na região Sul, a moda ficou dividida entre os conceitos 3 e 4, \textbf{cada} um deles atribuídos a três cursos ( 30,0 \% para \textbf{cada} um ), da mesma forma que na região Centro-Oeste, cujos dois cursos participantes obtiveram conceitos 2 e 3, sendo um curso para \textbf{cada} um ( 50 \% \textbf{cada} ).
Na região Norte, não houve cursos participantes.
O conceito 4 foi o segundo mais frequente em nível nacional ( 23,5 \%, correspondendo a oito cursos ), e o conceito 5, o terceiro ( 8,8 \%, correspondendo a três cursos ).
Houve, ainda, um curso ( 2,9 \% ) que recebeu conceito 1 e outro ( 2,9 \% ) que recebeu conceito 2.
Dos 34 cursos de Letras-Português, cinco ( 14,7 \% ) ficaram Sem Conceito ( SC ).
A região Norte não teve nenhum curso participante.
A região Nordeste contou com dois cursos participantes ou 5,9 \% do total nacional, que obtiveram, ambos, o conceito 3, como já comentado.
Dos 20 cursos participantes da região Sudeste ( 58,8 \% do total nacional ), o conceito 3 foi o conceito modal, como já anteriormente assinalado, tendo sido obtido por dez cursos ( 50,0 \% ).
O conceito 1 foi atribuído a um curso ( 5,0 \% ).
Receberam os conceitos 4 e 5, respectivamente, cinco cursos ( 25,0 \% ) e um curso ( 5,0 \% ).
Nessa região, três cursos ( 15,0 \% ) ficaram Sem Conceito ( SC ).
Os dez cursos da região Sul corresponderam a 29,4 \% do total nacional.
Como já indicado, a moda ficou dividida entre os conceitos 3 e 4, atribuídos a 3 cursos ( 30,0 \% ), \textbf{cada} um.
Dois cursos ( 20,0 \% ) obtiveram conceito 5, e outros dois cursos ( 20,0 \% ) ficaram Sem Conceito ( SC ) na região Sul.
Já dos dois cursos participantes na região Centro-Oeste ( 5,9 \% do total nacional ), um recebeu conceito 2, e o outro recebeu conceito 3, como já destacado.
No Gráfico 5.1b, apresenta -se a distribuição cumulativa do Conceito Enade dos cursos de Bacharelado segundo a Grande Região.
Em um gráfico de distribuição cumulativa, em \textbf{cada} valor de abscissa ( eixo x ), o valor apresentado no eixo das ordenadas ( eixo y ) é o valor acumulado da variável.
Por exemplo, as regiões Centro-Oeste ( linha preta ) e Nordeste ( linha verde ) apresentam 100 \% dos cursos com conceito 3 ou abaixo deste valor, já que as poligonais cumulativas alcançam 100 \% no conceito 3.
A região Sul ( linha roxa ), por outro lado, apresenta um pouco mais de 37 \% dos cursos com conceito 3, pois nos conceitos 1 e 2, a poligonal cumulativa indica que não houve cursos, tendo acrescido a porcentagem apenas quando se considera o conceito 3.
Em linhas gerais, em um gráfico de distribuição cumulativa, poligonais mais à esquerda representam Grandes Regiões com pior distribuição, e poligonais mais à direita, Grandes Regiões com melhores conceitos.
A região Centro-Oeste ( que só apresenta cursos com conceito 2 e 3 ) e a região Nordeste ( que só apresenta cursos com conceito 3 ) apresentam os cursos com conceitos mais baixos, ao passo que as regiões Sul ( linha roxa ) e Sudeste ( linha vermelha ) apresentam as distribuições com valores maiores.
Não houve cursos em IES situadas na região Norte ( linha azul ). 5.2 CONCEITOS POR CATEGORIA ADMINISTRATIVA, POR MODALIDADE DE ENSINO E GRANDE REGIÃO Na Tabela 5.2a, apresenta -se a distribuição dos cursos de Licenciatura de Letras-Português participantes do Enade/2017, por Categoria Administrativa e Modalidade de ensino, de acordo com os conceitos por eles alcançados, segundo as Grandes Regiões brasileiras.
Dos 318 cursos participantes, 201 ( 63,2 \% ) eram ministrados em Instituições Públicas, e 117 ( 36,8 \% ), em Instituições Privadas.
Quanto à Modalidade de ensino, 86,2 \% dos cursos eram presenciais, e 13,8 \%, a distância.
De acordo com as informações da Tabela 5.2a, em termos nacionais, entre Instituições Públicas e Privadas, observa -se que, dos 20 cursos avaliados com conceito 5, 15 eram oferecidos em IES Públicas, e cinco, em IES Privadas.
Dos 201 cursos participantes de IES Públicas, o conceito 3 foi o valor modal, atribuído a 65 cursos ( 32,3 \% ).
Dentre os demais cursos de Instituições Públicas participantes, 18 cursos ( 9,0 \% ) foram avaliados com conceito 1, 53 cursos ( 26,4 \% ) obtiveram conceito 2, 48 cursos ( 23,9 \% ) receberam conceito 4 e, como já comentado, 15 cursos ( 7,5 \% ) receberam conceito 5.
Nessa categoria, dois cursos ( 1,0 \% ) ficaram Sem Conceito ( SC ).
Na rede privada, o conceito modal foi 3, atribuído a 43 cursos, o correspondente a 36,8 \% dos 117 cursos da categoria.
Dentre os demais cursos participantes, seis ( 5,1 \% ) receberam conceito 1, e 20 ( 17,1 \% ), conceito 2.
O conceito 4 foi atribuído a 30 cursos ( 25,6 \%, ), e o conceito 5, a cinco cursos ( 4,3 \% ).
Nessa Categoria Administrativa, 13 cursos ( 11,1 \% ) ficaram Sem Conceito ( SC ).
Dos 86,2 \% dos cursos participantes da modalidade Presencial ( 274 ), o conceito 3 foi o valor modal, atribuído a 32,8 \% dos cursos ( 90 ).
Dentre os demais cursos nessa modalidade de ensino, 26,3 \% ( 72 ) foram avaliados com conceito 4, 20,8 \% ( 57 ) obtiveram conceito 2, 7,3 \% ( 20 ) receberam conceito 5, e 7,3 \% ( 20 ) receberam conceito 1.
Nessa modalidade, 5,5 \% dos cursos ( 15 ) ficaram Sem Conceito ( SC ).
Dentre os 44 cursos a Distância, o conceito modal foi 3, atribuído a 40,9 \% dos cursos da modalidade ( 18 ).
Dentre os demais cursos participantes, 13,6 \% ( seis cursos ) receberam conceito 4, 36,4 \% ( 16 cursos ) receberam conceito 2, e 9,1 \% ( quatro cursos ), conceito 1.
Nessa modalidade, nenhum curso ficou Sem Conceito ( SC ) ou recebeu conceito 5.
Na análise por região, observa -se que, na região Norte, as Instituições Privadas participaram com oito cursos ( 21,1 \% do total regional ), dos quais dois ( 25,0 \% ) ficaram Sem Conceito ( SC ), e nenhum recebeu os conceitos 4 ou 5.
O conceito modal para as Instituições Privadas na região foi 2, atribuído a três cursos, correspondendo a 37,5 \% dessa categoria na região.
Dois cursos ( 25,0 \% ) foram avaliados com conceito 3, e um curso ( 12,5 \% ) recebeu conceito 1.
As Instituições Públicas participaram com 30 cursos na região Norte ( 78,9 \% do total regional ).
A moda ficou dividida entre os conceitos 2 e 3, atribuídos, \textbf{cada} um, a dez cursos ( 33,3 \% ).
Os conceitos 1 e 4 também foram atribuídos ao mesmo número de cursos, no caso, a cinco cursos ( 16,7 \% ) \textbf{cada} um.
Nenhum curso ficou Sem Conceito ( SC ) ou com conceito 5 nessa combinação de categoria e região.
Na região Norte, 94,7 \% dos cursos oferecidos eram presenciais ( 36 cursos ).
Desses, 12 ( 33,3 \% ) alcançaram o conceito modal 3.
Outros seis ( 16,7 \% ) cursos obtiveram conceito 1, 11 ( 30,6 \% ) obtiveram conceito 2, e cinco cursos ( 13,9 \% ), conceito 4.
Dois cursos ( 5,6 \% ) ficaram Sem Conceito ( SC ), e nenhum recebeu conceito 5.
A região Norte teve os seus dois únicos cursos a Distância recebendo conceito 2.
Na região Nordeste, a rede privada concentrou 16 dos 113 cursos participantes, o equivalente a 14,2 \% do total da região.
Dentre os cursos oferecidos por IES Privadas no Nordeste, nove cursos ( 56,3 \% ) foram avaliados com conceito 3, o conceito modal.
Três cursos ( 18,8 \% ) obtiveram conceito 2, e outros três ( 18,8 \% ) obtiveram conceito 4.
Um curso ( 6,3 \% ) ficou Sem Conceito ( SC ).
Nessa combinação de categoria e região, nenhum curso recebeu os conceitos 1 ou 5.
As Instituições Públicas da região Nordeste participaram com 97 cursos ( 85,8 \% do total da região ).
Desses, 38 ( 39,2 \% ) obtiveram conceito 3, o valor modal.
Outros 11 cursos ( 11,3 \% ) obtiveram conceito 1, 29 cursos ( 29,9 \% ) obtiveram conceito 2.
O conceito 4 foi atribuído a 16 cursos ( 16,5 \% ), e o conceito 5, a dois cursos ( 2,1 \% ).
Um curso ( 1,0 \% ) ficou Sem Conceito ( SC ) nessa região.
A maioria dos cursos na região Nordeste foram presenciais : 85,8 \% ( 97 cursos ).
Para esses, a moda foi o conceito 3, com 43,3 \% dos cursos ( 42 ).
Os demais cursos presenciais da região foram alocados aos conceitos 1 ( 9,3 \%, nove cursos ), 2 ( 25,8 \%, 25 cursos ), 4 ( 17,5 \%, 17 cursos ) e 5 ( 2,1 \%, dois cursos ).
Dois cursos presenciais ( 2,1 \% ) ficaram Sem Conceito ( SC ).
Os cursos a Distância ( 14,2 \%, 16 cursos ) tiveram o conceito 2 como o valor modal, atribuído a 43,8 \% dos cursos ( sete cursos ).
Nenhum curso a Distância da região ficou Sem Conceito ( SC ) ou recebeu conceito 5.
O conceito 1 foi atribuído a 12,5 \% dos cursos ( dois ), o conceito 3, a 31,3 \% dos cursos ( cinco ), e o conceito 4, a 12,5 \% dos cursos ( dois ).
Na região Sudeste, a proporção de cursos da rede privada, 58,9 \%, foi mais elevada do que a registrada nas demais regiões brasileiras, correspondendo a 53 dos 90 cursos participantes.
Nessa categoria e região, o conceito modal foi 4, atribuído a 15 cursos ( 28,3 \% ).
Sete cursos ( 13,2 \% ) ficaram Sem Conceito ( SC ).
Os demais foram avaliados com conceito 1 ( dois cursos, 3,8 \% ), conceito 2 ( 11 cursos, 20,8 \% ), conceito 3 ( 14 cursos, 26,4 \% ) e conceito 5 ( quatro cursos, 7,5 \% ).
Nos 37 cursos ( 41,1 \% ) oferecidos por Instituições Públicas na região Sudeste, o conceito modal foi 4, atribuído a 15 cursos ( 40,5 \% ).
Os demais cursos receberam os conceitos 2 ( seis cursos, 16,2 \% ), 3 ( sete cursos, 18,9 \% ) e 5 ( oito cursos, 21,6 \% ).
Na região Sudeste, nessa categoria, nenhum curso recebeu conceito 1 e apenas um curso ( 2,7 \% ) ficou Sem Conceito ( SC ).
Também na região Sudeste, a maioria dos cursos foram presenciais : 85,6 \% ( 77 cursos ).
Para esses, a moda foi o conceito 4, com 35,1 \% dos cursos ( 27 ).
Os demais cursos presenciais da região foram alocados aos conceitos 1 ( 2,6 \%, dois cursos ), 2 ( 15,6 \%, 12 cursos ), 3 ( 20,8 \%, 16 cursos ) e 5 ( 15,6 \%, 12 cursos ).
Oito cursos presenciais ( 10,4 \% ) ficaram Sem Conceito ( SC ).
Os 13 cursos a Distância, 14,4 \% na região, tiveram a moda dividida entre os conceitos 2 e 3 ( 38,5 \%, equivalentes a 5 cursos, para \textbf{cada} ).
Os demais cursos a distância da região ( três cursos ) foram alocados ao conceito 4 ( 23,1 \% dos cursos ).
Nenhum curso a Distância da região foi alocado ao conceito 1 ou 5, assim como nenhum ficou Sem Conceito ( SC ).
As Instituições Privadas concentraram 26 dos 46 cursos participantes da região Sul, 56,5 \% do total regional.
Desses, 11 ( 42,3 \% ) obtiveram conceito 3, o conceito modal.
Nessa combinação de Categoria Administrativa e Grande Região, o conceito 2 foi atribuído a três cursos ( 11,5 \% ).
Os conceitos 4 e 5 foram atribuídos, respectivamente, a dez cursos ( 38,5 \% ) e a um curso ( 3,8 \% ).
Nenhum curso recebeu conceito 1, e um curso ( 3,8 \% ) ficou Sem Conceito ( SC ).
As Instituições Públicas na região Sul participaram com 20 cursos, 43,5 \% dos cursos da região.
O conceito modal foi 4, atribuído a dez cursos ( 50,0 \% ).
Aos demais cursos foram atribuídos os conceitos 3 ( quatro cursos, 20,0 \% ) e 2 ( dois cursos, 10,0 \% ).
O conceito 5 foi atribuído a quatro cursos ( 20,0 \% ).
Nenhum curso da região Sul oferecido por IES Pública ficou Sem Conceito ( SC ) ou recebeu conceito 1.
Também na região Sul, a maioria dos cursos foram presenciais, correspondendo a 82,6 \% ( 38 cursos ), em oposição aos 17,4 \% dos cursos a Distância ( oito cursos ).
Para os cursos presenciais, a moda foi o conceito 4, com 52,6 \% ( 20 ) dos cursos.
Os demais cursos presenciais da região foram alocados aos conceitos 2 ( 7,9 \%, três cursos ), 3 ( 23,7 \%, nove cursos ) e 5 ( 13,2 \%, cinco cursos ).
Nenhum curso presencial foi alocado ao conceito 1, e apenas um curso ( 2,6 \% ) ficou Sem Conceito ( SC ).
Para os oito cursos a Distância, a moda foi o conceito 3 ( 75,0 \%, seis cursos ).
Outros 25,0 \% dos cursos a Distância da região foram alocados ao conceito 2 ( dois cursos ).
Nenhum curso a Distância da região foi alocado aos demais conceitos ou ficou Sem Conceito ( SC ).
Na região Centro-Oeste, 14 dos 31 cursos participantes eram de Instituições Privadas ( 45,2 \% em termos regionais ).
Desses, sete ( 50,0 \% ) receberam conceito 3, o conceito modal.
Dos demais, três ( 21,4 \% ) receberam conceito 1, e dois ( 14,3 \% ), conceito 4.
Dois cursos ( 14,3 \% ) ficaram Sem Conceito ( SC ), e nenhum curso recebeu os conceitos 2 ou 5.
Quanto aos 17 cursos oferecidos por Instituições Públicas na região Centro-Oeste ( 54,8 \% do total regional ), a moda ficou dividida entre os conceitos 2 e 3, \textbf{cada} um com seis cursos ( 35,3 \% \textbf{cada} ).
Dois cursos ( 11,8 \% ) apresentaram conceito 1, outros dois ( 11,8 \% ) receberam conceito 4, e um curso ( 5,9 \% ) recebeu conceito 5.
Nenhum curso ficou Sem Conceito ( SC ).
Nessa região, 83,9 \% dos cursos foram presenciais ( 26 cursos ), e 16,1 \%, a Distância ( cinco cursos ).
Para os cursos presenciais, a moda foi o conceito 3, com 42,3 \% dos cursos ( 11 ).
Os demais foram alocados aos conceitos 1 ( 11,5 \%, três cursos ), 2 ( 23,1 \%, seis cursos ), 4 ( 11,5 \%, três cursos ), e 5 ( 3,8 \%, um curso ).
Dois cursos presenciais ( 7,7 \% ) ficaram Sem Conceito ( SC ).
Entre os cinco cursos a Distância participantes a moda ficou dividida igualmente nos conceitos 1 e 3 ( 40,0 \%, dois cursos para \textbf{cada} um ) e um curso ( 20,0 \% ) recebeu o conceito 4.
No Gráfico 5.2a, é apresentada a distribuição cumulativa do Conceito Enade dos cursos de Licenciatura, segundo a Categoria Administrativa da IES.
Os cursos em IES Públicas ( linha azul ) apresentam uma poligonal bem semelhante da poligonal dos cursos em IES Privadas ( linha verde ), onde a porcentagem de cursos com conceitos mais baixos é maior para cursos em IES Públicas, mas a porcentagem de cursos com conceito máximo é maior para cursos em IES Públicas.
Na Tabela 5.2b, é apresentada a distribuição dos cursos de Bacharelado participantes do Enade/2017 de Letras-Português, por Categoria Administrativa e Modalidade de ensino, de acordo com os conceitos por eles alcançados, segundo as Grandes Regiões brasileiras.
Dos 34 cursos participantes, 26 ( 76,5 \% ) eram ministrados em Instituições Públicas, e oito ( 23,5 \% ), em Instituições Privadas.
Quanto à Modalidade de ensino, todos os 34 cursos eram presenciais.
A moda entre eles foi o conceito 3, atribuído a 16 cursos ( 47,1 \% ).
Os conceitos 1 e 2 foram atribuídos a um curso \textbf{cada} ( 2,9 \% para \textbf{cada} um ), o conceito 4, a oito cursos ( 23,5 \% ), e o conceito 5, a três cursos ( 8,8 \% ).
Cinco cursos ( 14,7 \% ) ficaram Sem Conceito ( SC ).
Como já mencionado anteriormente, não houve participação de cursos na Modalidade a Distância.
De acordo com as informações da Tabela 5.2b, em termos nacionais, dentre Instituições Públicas e Privadas, observa -se que, todos os três cursos avaliados com conceito 5 eram oferecidos em IES Públicas.
Dos 26 cursos participantes de IES Públicas, o conceito 3 foi o valor modal, atribuído a 12 cursos ( 46,2 \% ).
Dentre os demais cursos de Instituições Públicas participantes, um curso ( 3,8 \% ) foi avaliado com conceito 1, um curso ( 3,8 \% ), com conceito 2, seis cursos ( 23,1 \% ) foram avaliados com conceito 4, e três cursos ( 11,5 \% ), com conceito 5.
Nessa categoria, três cursos ( 11,5 \% ) ficaram Sem Conceito ( SC ).
Na rede privada, o conceito modal foi 3, atribuído a quatro cursos, o correspondente a 50,0 \% dos oito cursos da categoria.
Dos demais cursos participantes, dois ( 25,0 \% ) receberam conceito 4, e outros dois ( 25,0 \% ) ficaram Sem Conceito ( SC ) nessa Categoria Administrativa.
Na análise por região, observa -se que, na região Norte, não houve cursos inscritos.
Na região Nordeste, a rede pública participou com os dois únicos cursos inscritos, que receberam conceito 3.
Na região Sudeste, a proporção de cursos da rede privada, 40,0 \%, foi mais elevada do que a registrada nas demais regiões brasileiras, correspondendo a oito dos 20 cursos participantes.
Nessa categoria e região, o conceito modal foi 3, atribuído a quatro cursos ( 50,0 \% ).
Dois cursos ( 25,0 \% ) ficaram Sem Conceito ( SC ), e outros dois cursos ( 25,0 \% ) foram avaliados com conceito 4.
Nos 12 cursos oferecidos por Instituições Públicas na região Sudeste ( 60,0 \% em termos regionais ), o conceito modal foi 3, atribuído a seis cursos ( 50,0 \% ).
Os demais cursos receberam os conceitos 1 ( um curso, 8,3 \% ), 4 ( três cursos, 25,0 \% ) e 5 ( um curso, 8,3 \% ).
Na região Sudeste, nessa categoria, um curso ( 8,3 \% ) ficou Sem Conceito ( SC ).
Nenhum curso recebeu conceito 2.
As Instituições Públicas concentraram todos os dez cursos participantes da região Sul, 100,0 \% do total regional.
Desses, três cursos obtiveram conceito 3 ( 30,0 \% ), e outros três ( 30,0 \% ), conceito 4.
Dois cursos ( 20,0 \% ) obtiveram conceito 5, e dois ( 20,0 \% ) ficaram Sem Conceito ( SC ).
Nenhum curso foi alocado aos conceitos 1 ou 2.
As Instituições Privadas na região Sul não participaram com nenhum curso.
Na região Centro-Oeste, apenas dois cursos participaram, ambos de Instituições Públicas.
Um recebeu conceito 2, e outro, conceito 3.
No Gráfico 5.2b, é apresentada a distribuição cumulativa do Conceito Enade dos cursos de Bacharelado, segundo a Categoria Administrativa da IES.
Os cursos em IES Públicas ( linha azul ) apresentam uma poligonal mais à direita do que os em IES Privadas ( linha verde ) e, consequentemente, uma distribuição de Conceitos Enade com valores maiores, já que o acréscimo do percentual nos conceitos mais altos é maior para cursos de IES Públicas do que para os de IES Privadas. 5.3 CONCEITOS POR ORGANIZAÇÃO ACADÊMICA E GRANDE REGIÃO Na Tabela 5.3a, encontra -se a distribuição dos conceitos atribuídos aos cursos de Licenciatura participantes do Enade/2017 na Área de Letras-Português, por Organização Acadêmica, segundo as Grandes Regiões brasileiras.
Dos 318 cursos de Letras-Português participantes, 236 eram oferecidos em Universidades, 21 em Centros Universitários, 49 em Faculdades e 12 em CEFET/IFET.
Essa distribuição corresponde a, respectivamente, 74,2 \%, 6,6 \%, 15,4 \% e 3,8 \% dos cursos.
De acordo com os dados apresentados, dos 20 cursos avaliados com conceito 5, 15 eram vinculados a Universidades.
Os cursos vinculados a esse tipo de Organização Acadêmica tiveram o conceito 3 como conceito modal, atribuído a 82 cursos ( 34,7 \% ).
Os demais cursos vinculados a Universidades receberam os conceitos 1 ( 19 cursos, 8,1 \% ), 2 ( 54 cursos, 22,9 \% ), 4 ( 57 cursos, 24,2 \% ) e 5 ( 15 cursos, 6,4 \%, como já mencionado ).
Ficaram Sem Conceito ( SC ) nove cursos ( 3,8 \% ).
Dentre os cursos em Centros Universitários, o conceito modal foi 4, atribuído a dez cursos ( 47,6 \% ).
Um curso vinculado a esse tipo de Organização Acadêmica ( 4,8 \% ) ficou Sem Conceito ( SC ).
Cinco cursos ( 23,8 \% ) receberam os conceitos 2 e 3, \textbf{cada} um.
Nenhum curso recebeu os conceitos 1 ou 5.
Dos 49 cursos mantidos por Faculdades, 15 ( 30,6 \% ) receberam conceito 3, o conceito modal.
O segundo conceito mais frequentemente atribuído aos cursos oferecidos por essa categoria de IES foi o conceito 2, recebido por 12 cursos ( 24,5 \% ).
Na sequência, dez cursos ( 20,4 \% ) obtiveram conceito 4, e dois cursos ( 4,1 \% ), conceito 5.
Cinco cursos ( 10,2 \% ) ficaram Sem Conceito ( SC ) e a outros cinco ( 10,2 \% ) foram atribuídos o conceito 1.
Dos 12 cursos mantidos por CEFET/IFET, seis ( 50,0 \% ) receberam conceito 3, o valor modal.
O segundo conceito mais frequente atribuído aos cursos oferecidos por essa categoria de IES foi o conceito 5, recebido por três cursos ( 25,0 \% ).
Os conceitos 2 e 4 foram atribuídos, respectivamente, a dois cursos ( 16,7 \% ) e a um curso ( 8,3 \% ).
Nenhum curso ficou Sem Conceito ( SC ) ou recebeu o conceito 1.
Considerando -se separadamente as regiões brasileiras, verifica -se que, na região Norte, as Universidades participaram com 30 ( 78,9 \% ) dos 38 cursos participantes.
Desses, 11 cursos ( 36,7 \% ) receberam conceito 2, conceito modal.
Os demais cursos receberam os conceitos 1 ( cinco cursos, 16,7 \% ), 3 ( dez cursos, 33,3 \% ) e 4 ( quatro cursos, 13,3 \% ).
Nenhum curso ficou Sem Conceito ( SC ) ou recebeu conceito 5 nessa combinação de categoria e região.
Os Centros Universitários da região Norte participaram com apenas um curso da Área de Letras-Português ( Licenciatura, o qual ficou Sem Conceito ( SC ).
As Faculdades participaram com cinco cursos na região Norte ( 13,2 \% em termos regionais ) que receberam os conceitos 2 ( dois cursos, 40,0 \%, o conceito modal ), 1 e 3 ( um curso, 20,0 \% \textbf{cada} um ).
Um curso ( 20,0 \% ) ficou Sem Conceito ( SC ).
Nenhum curso de Faculdades da região obteve os conceitos 4 ou 5.
Os CEFET/IFET participaram com dois cursos. \textbf{Cada} um deles recebeu os conceitos 3 e 4.
Na região Nordeste, as Universidades contaram com 97 ( 85,8 \% ) dos 113 cursos oferecidos, a maior proporção entre as regiões dessa categoria.
O conceito modal 3 foi obtido por 38 cursos ( 39,2 \% ).
Dois cursos ( 2,1 \% ) ficaram Sem Conceito ( SC ).
Outros 11 cursos ( 11,3 \% ) receberam conceito 1, 28 ( 28,9 \% ) receberam conceito 2, 16 ( 16,5 \% ) receberam conceito 4, e dois ( 2,1 \% ), conceito 5.
Os Centros Universitários contaram com três cursos participantes na região Nordeste, 2,7 \% em termos regionais.
Desses, dois cursos ( 66,7 \% ) obtiveram conceito modal 4, e um curso ( 33,3 \% ) obteve conceito 3.
As Faculdades foram representadas por nove cursos na região Nordeste, 8,0 \% em termos regionais.
O conceito modal 3 foi atribuído a cinco cursos ( 55,6 \% ), o conceito 2, a três cursos ( 33,3 \% ), e o conceito 4, a um curso ( 11,1 \% ).
Os CEFET/IFET participaram com quatro cursos, 3,5 \% dos cursos na região, dos quais três ( 75,0 \% ) receberam o conceito 3 e um curso ( 25,0 \% ) o conceito 2.
Na região Sudeste, as Universidades concentraram 54 ( 60,0 \% ) dos 90 cursos de Letras-Português ( Licenciatura ) da região.
Dos cursos oferecidos em Universidades, nessa região, o conceito modal foi 4, atribuído a 17 cursos ( 31,5 \% ).
Os demais cursos receberam os conceitos 1 ( um curso, 1,9 \% ), 2 ( oito cursos, 14,8 \% ), 3 ( 14 cursos, 25,9 \% ) e 5 ( oito cursos, 14,8 \% ).
Seis cursos ( 11,1 \% ) ficaram Sem Conceito ( SC ).
Os Centros Universitários participaram com 11 cursos na região Sudeste, equivalentes a 12,2 \% em termos regionais ( a maior proporção dessa categoria entre as regiões ), dos quais sete ( 63,6 \% ) obtiveram o conceito modal 4.
Os conceitos 2 e 3 foram atribuídos, \textbf{cada} um, a dois cursos ( 18,2 \% \textbf{cada} ).
Nenhum curso ficou Sem Conceito ( SC ) nessa combinação de organização acadêmica e região, tampouco recebeu os conceitos 1 ou 5.
Já as Faculdades foram representadas por 20 cursos na região Sudeste, 22,2 \% em termos regionais, dos quais dois ( 10,0 \% ) ficaram Sem Conceito ( SC ).
A moda dividiu -se entre os conceitos 2 e 4, atribuídos a seis cursos, \textbf{cada} um ( 30,0 \% para \textbf{cada} conceito ).
Aos demais cursos foram atribuídos os conceitos 1 ( um curso, 5,0 \% ), 3 ( quatro cursos, 20,0 \% ) e 5 ( um curso, 5,0 \% ).
Os CEFET/IFET participaram com cinco cursos, 5,6 \% em termos regionais, os quais receberam conceitos 5 ( três cursos, 60,0 \%, o conceito modal ), 2 e 3 ( um curso, 20,0 \% \textbf{cada} um ).
Dos 46 cursos da região Sul, 36 ( 78,3 \% ) eram de Universidades, para os quais o conceito modal foi 4, atribuído a 17 cursos ( 47,2 \% ).
Nesse tipo de organização, um curso ( 2,8 \% ) ficou Sem Conceito ( SC ), e nenhum curso recebeu conceito 1.
Os demais cursos receberam os conceitos 2 ( dois cursos, 5,6 \% ), 3 ( 12 cursos, 33,3 \% ) e 5 ( quatro cursos, 11,1 \% ).
Dos cinco cursos ( 10,9 \% ) participantes de Centros Universitários da região Sul, nenhum ficou Sem Conceito ( SC ) ou recebeu os conceitos 1 ou 5.
O conceito modal, 2, foi atribuído a três cursos ( 60,0 \% ).
Os conceitos restantes, 3 e 4, foram atribuídos a um curso ( 20,0 \% ), \textbf{cada} um.
Quanto aos também cinco cursos ( 10,9 \% ) vinculados a Faculdades na região Sul, a moda ficou dividida entre os conceitos 3 e 4, atribuídos a dois cursos ( 40,0 \% ), \textbf{cada} um.
O conceito 5 foi atribuído a um curso ( 20,0 \% ).
Nenhum curso ficou Sem Conceito ( SC ) ou recebeu os conceitos 1 ou 2.
Nessa região, não foram avaliados cursos em CEFET/IFET.
Na região Centro-Oeste, 19 ( 61,3 \% ) dos 31 cursos eram de Universidades.
Nesse tipo de organização, nenhum curso ficou Sem Conceito ( SC ).
Os conceitos 1, 2, 4 e 5 foram, respectivamente, atribuídos a dois ( 10,5 \% ), cinco ( 26,3 \% ), três ( 15,8 \% ) cursos e um curso ( 5,3 \% ).
O conceito 3, conceito modal, foi atribuído a oito cursos ( 42,1 \% ).
Os Centros Universitários da região Centro-Oeste contaram com apenas um curso, que obteve conceito 3.
Já as Faculdades contaram com dez cursos, 32,3 \% dos cursos na região ( a maior proporção dessa categoria entre as regiões ), dos quais dois ( 20,0 \% ) ficaram Sem Conceito ( SC ).
Os conceitos 1 e 3, com três cursos \textbf{cada} um ( 30,0 \% para \textbf{cada} ) foram os conceitos modais.
Os conceitos 2 e 4 foram atribuídos, \textbf{cada} um, a um curso ( 10,0 \% para \textbf{cada} ).
Nenhum curso recebeu conceito 5.
Apenas um curso foi avaliado em CEFET/IFET, tendo obtido conceito 3.
No Gráfico 5.3a, é apresentada a distribuição cumulativa do Conceito Enade dos cursos de Licenciatura, segundo a Organização Acadêmica da IES.
Os cursos em CEFET/IFET ( linha roxa ) e em Centros Universitários ( linha verde ) apresentam uma poligonal mais à direita do que as demais, denotando melhor distribuição de Conceitos Enade, mesmo com a poligonal cumulativa dos cursos de Centros Universitários ( linha verde ) não apresentando cursos com conceito máximo.
No outro extremo, encontram -se os cursos oferecidos em Faculdades ( linha vermelha ) com a poligonal mais à esquerda e, consequentemente, uma distribuição de Conceitos Enade com menores valores.
IES de Universidades ( linha azul ) apresentam uma situação intermediária.
Na Tabela 5.3b, encontra -se a distribuição dos conceitos atribuídos aos cursos de Bacharelado participantes do Enade/2017 na Área de Letras-Português, por Organização Acadêmica, segundo as Grandes Regiões brasileiras.
Dos 34 cursos de Letras-Português participantes, 31 eram oferecidos em Universidades, nenhum em Centros Universitários, dois em Faculdades e um em CEFET/IFET.
Essa distribuição corresponde a, respectivamente, 91,2 \%, 0,0 \%, 5,9 \% e 2,9 \% dos cursos.
De acordo com os dados apresentados, todos os três cursos avaliados com conceito 5 eram vinculados a Universidades.
Os cursos vinculados a esse tipo de Organização Acadêmica tiveram o conceito 3 como conceito modal, atribuído a 14 cursos ( 45,2 \% ).
Os demais cursos vinculados a Universidades receberam os conceitos 1 e 2 ( um curso, 3,2 \% \textbf{cada} um ), 4 ( oito cursos, 25,8 \% ) e 5 ( três cursos, 9,7 \%, como já mencionado ).
Ficaram Sem Conceito ( SC ) quatro cursos ( 12,9 \% ).
Como já mencionado, não houve representantes dos Centros Universitários.
Dos dois cursos mantidos por Faculdades, um ficou Sem Conceito ( SC ), e o outro obteve conceito 3.
O único curso mantido por CEFET/IFET recebeu conceito 3.
Considerando -se separadamente as regiões brasileiras, verifica -se que, na região Norte, como já mencionado, não houve cursos participantes.
Na região Nordeste, as Universidades participaram com apenas dois cursos da Área de Letras-Português, que obtiveram conceito 3.
Centros Universitários, Faculdades e CEFET/IFET não contaram com cursos participantes na região Nordeste.
Na região Sudeste, as Universidades concentraram 17 ( 85,0 \% ) dos 20 cursos de Letras-Português da região.
Dos cursos em Universidades, nessa região, o conceito modal foi 3, atribuído a oito cursos ( 47,1 \% ).
Os demais cursos receberam os conceitos 1 ( um curso, 5,9 \% ), 4 ( cinco cursos, 29,4 \% ) e 5 ( um curso, 5,9 \% ).
Nenhum curso recebeu conceito 2, e dois cursos ( 11,8 \% ) ficaram Sem Conceito ( SC ).
Os Centros Universitários não participaram com cursos na região Sudeste.
Já as Faculdades foram representadas por dois cursos na região, dos quais um ficou Sem Conceito ( SC ), e o outro recebeu conceito 3.
Houve participação de apenas um curso em CEFET/IFET nessa região, que obteve conceito 3.
Dos dez cursos da região Sul, todos eram de Universidades, cuja moda ficou dividida entre os conceitos 3 e 4, atribuídos a três cursos, \textbf{cada} um ( 30,0 \% \textbf{cada} ).
Nesse tipo de organização, dois cursos ( 20,0 \% ) ficaram Sem Conceito ( SC ) e outros dois cursos ( 20,0 \% ) receberam conceito 5.
Não houve cursos participantes em Faculdades, Centros Universitários ou CEFET/IFET.
Na região Centro-Oeste, apenas dois cursos participaram, ambos de nas Universidades, tendo obtido os conceitos 2 e 3, atribuídos, \textbf{cada} um, a um curso.
Centros Universitários, Faculdades e CEFET/IFET não tiveram cursos representantes na região Centro-Oeste.
No Gráfico 5.3b, é apresentada a distribuição do Conceito Enade dos cursos de Bacharelado, segundo a Organização Acadêmica da IES.
Como já mencionado, não houve cursos representando Centros Universitários.
Os cursos em Universidades ( linha azul ) apresentam uma poligonal mais à direita do que as demais, denotando uma melhor distribuição de Conceitos Enade.
No outro extremo, encontram -se os cursos oferecidos em Faculdades ( linha verde ) e em CEFET/IFET ( linha vermelha ), com as poligonais coincidindo mais à esquerda e, consequentemente, uma distribuição de Conceitos Enade com menores valores.

% matched lemmas: cada, estar
\end{textsample}
